\documentclass[11pt]{article}
\usepackage[margin=25mm]{geometry}
\usepackage{fontspec}
\setmainfont{Latin Modern Roman}
\usepackage{amsmath,amssymb,mathtools}
\usepackage{unicode-math}
\setmathfont{Latin Modern Math}
\usepackage{xurl,hyperref}
\hypersetup{colorlinks=true,urlcolor=blue}
\setcounter{secnumdepth}{0}
\setlength{\parindent}{0pt}
\setlength{\parskip}{5pt}
\setlength{\emergencystretch}{3em}
\newcommand{\sref}[2]{\hyperlink{source:#1:#2}{[S#2]}}
\newcommand{\cataloguescope}[1]{\paragraph{Recorded target.}#1}
\newcommand{\R}{\mathbb R}
\newcommand{\C}{\mathbb C}
\newcommand{\N}{\mathbb N}
\newcommand{\Z}{\mathbb Z}
\newcommand{\Q}{\mathbb Q}
\newcommand{\F}{\mathbb F}
\newcommand{\eps}{\varepsilon}
\begin{document}
% UnsolvedMath ID 30004805; source catalogue code OWR-8415343-008
\setcounter{section}{30004804}
\section{Kasner Singularities Beyond the Subcritical Regime}\label{30004805}
{\small UnsolvedMath ID 30004805 · Mathematical Physics}
\subsection{Definitions and mathematical statement}

\paragraph{Shared notation.}$\mathbb N=\{1,2,\ldots\}$; the standard number fields have their usual meanings. All additional conventions are stated locally.


A generalized Kasner background has spatial scale factors $t^{\widetilde q_I}$ as proper time $t$ approaches zero. The exponents are those of the vacuum or Einstein--scalar-field background in the report. Asymptotically velocity-term-dominated behavior means that spatial derivative terms become negligible relative to time evolution. The subcritical condition is the displayed strict inequality in the target.

\paragraph{Catalogue wording retained for traceability.}
For general sufficiently small perturbations, possibly without symmetry, of generalized Kasner solutions of the Einstein equations, what behavior occurs near the singularity when the sub-criticality condition $\max_{I<J,B}(\widetilde q_I+\widetilde q_J-\widetilde q_B)<1$ fails?

\paragraph{Scope and source qualification.}The failure of the strict inequality includes the borderline equality case. Oscillatory behavior is a heuristic possibility, rather than an added asserted conclusion.
\subsection{Short English statement}
Determine the singular behavior of small unrestricted perturbations of Kasner solutions outside the proven subcritical regime.
\subsection{Research attempt}
\paragraph{Model, date and investigated scope.}
OpenAI GPT-6 Astra (configured model \texttt{gpt-6-astra}), ultra reasoning; generated 11 October 2026. Derived the spatial-frame scaling behind the threshold and explicit borderline accumulation; no nonlinear beyond-subcritical asymptotics.

\paragraph{Sources and literature check.}
Read the supplied definition, catalogue wording, source qualifications and reference entries for 30004805. The historical source/context and inherited status are retained; no fresh literature or novelty certification is claimed. No external theorem is used without proof in this attempt; explicit model assumptions are distinguished from facts about the source model.

\paragraph{Where the subcritical power arises.} Let \(X_I\) be a time-independent spatial frame with \([X_I,X_J]=\sum_B c_{IJ}^B X_B\), and rescale it by Kasner powers \(e_I=t^{-q_I}X_I\). Since the powers depend only on time,
\[
 [e_I,e_J]=\sum_B c_{IJ}^B t^{q_B-q_I-q_J}e_B.
\]
Relative to the time-derivative scale \(t^{-1}\), a nonzero structure coefficient therefore has size \(t^{1-(q_I+q_J-q_B)}\). It decays when the exponent in parentheses is less than one, remains order one at equality, and grows when it exceeds one.
\paragraph{Borderline accumulation.} Even a coefficient of small amplitude cannot be dismissed merely for being small if its time integral is borderline: \(\int_t^1\varepsilon s^{-1}ds=\varepsilon\log(1/t)\) becomes order one at \(t=e^{-1/\varepsilon}\). For powers \(s^{-b}\), \(b>1\), accumulation is stronger, equal to \(\varepsilon(t^{1-b}-1)/(b-1)\). These exact integrations identify why extending a subcritical integrability estimate fails.
\paragraph{Scope audit.} The frame calculation does not prove that every potentially dangerous coefficient is nonzero, survives constraints or causes oscillations. Homogeneous commuting frames have all \(c_{IJ}^B=0\), showing that failure of the numerical condition alone cannot force this mechanism for every solution. The first gap is the nonlinear constrained evolution of the excited spatial-curvature modes, including the equality case, for generic small unrestricted perturbations. No new singularity classification is claimed.

\paragraph{Verification and adversarial audit.}
Checked the frame rescaling, equality versus supercritical powers and the commuting-frame countercheck; did not infer oscillation from failure of a sufficient condition. This is a GPT-6 Astra Ultra self-audit, not independent human verification. Prompt A drafting followed by prompt B scope/dependency/boundary self-audit. Checked the frame rescaling, equality versus supercritical powers and the commuting-frame countercheck; did not infer oscillation from failure of a sufficient condition. The remaining gap is: Control the constrained nonlinear evolution of genuinely excited curvature modes at and above the threshold, including possible cancellations. Unresolved status and zero original-target completion are retained.

Coordinator checked the rescaled frame commutator and logarithmic/power accumulation at the threshold, without asserting generic excitation of the mode. This additional check was model-assisted, not independent human verification.

\paragraph{Tools, process and contributions.}
Python 3 and local file-reading/writing tools for manuscript serialization and structural checks; mathematical calculations were symbolic reasoning, with no numerical experiment or theorem prover. Contribution prompt A was followed by prompt B. GPT-6 Astra with ultra reasoning drafted the research and performed its stated self-audit; any additional coordinator contribution is identified in the verification and audit. A GPT-6-based Codex coordinator assembled the manuscript and checked repository integration; its exact variant and reasoning configuration were not exposed. No human mathematical contribution, independent human proof check or formal proof certification is supplied. Complete derivations and the consulted source references are included here; private reasoning transcripts are not published.

\paragraph{Final status and first remaining gap.}
\textbf{UNRESOLVED}. Control the constrained nonlinear evolution of genuinely excited curvature modes at and above the threshold, including possible cancellations.

\section{Completion Estimate}
\noindent\textbf{COMPLETION: 0\%}
The proved calculation is an elementary auxiliary result or a test of a restricted strategy; no new portion of the full source target is established.

\subsection{Sources}
\begin{itemize}
\item[\hypertarget{source:30004805:1}{S1}] ULAM.AI and the UnsolvedMath Contributors, UnsolvedMath, record 30004805 (OWR-8415343-008). Catalogue curation under CC BY 4.0; definitions and English exposition added locally.
\url{https://huggingface.co/datasets/ulamai/UnsolvedMath}.
\item[\hypertarget{source:30004805:2}{S2}] Mathematical Aspects of General Relativity (2021), Oberwolfach Reports, DOI 10.4171/owr/2021/40, pp. 2157--2267. Original report and mathematical context.
\url{https://doi.org/10.4171/owr/2021/40}.
\end{itemize}
\label{end:30004805}
\end{document}
